\documentclass[11pt,a4paper]{article}

% --- Packages (imports from DistCat v4 main.tex) ---
\usepackage{amsmath,amssymb,amsthm,amsfonts}
\usepackage{physics}
\usepackage{geometry}
\geometry{margin=1in}
\usepackage{graphicx}
\usepackage{hyperref}
\usepackage{bm}
\hypersetup{hidelinks}

% Theorem-style environments (consistent with main.tex)
\newtheorem{theorem}{Theorem}[section]
\newtheorem{definition}[theorem]{Definition}
\newtheorem{proposition}[theorem]{Proposition}
\newtheorem{lemma}[theorem]{Lemma}
\newtheorem{corollary}[theorem]{Corollary}


\subsection{Example 1: D1-grav-waves --- AQFT-based Gravitational-wave
Source Discrimination in PTA}\label{sec:example1-grav-waves}

This section instantiates the DistCat framework for pulsar timing array
(PTA) gravitational-wave detection and source discrimination. We
establish that the Hellings--Downs (HD) correlation function is a
natural transformation in DistCat (Proposition~\ref{prop:hd-nat}), and
prove that distinct GW source classes are distinguishable by the total
variation distance between their algebraic states (Theorem~\ref{thm:source-discrim}),
yielding a provable discrimination threshold. Furthermore, we establish
the structural difference in anisotropy between supermassive black hole
binaries (SMBHBs, $\mathcal{O}(1\text{--}10\%)$) and cosmological
backgrounds ($\mathcal{O}(10^{-5})$) as a DistCat invariant
(Theorem~\ref{thm:isotropy-scaling}).

\textbf{Cross-reference}: This example corresponds to row 1 of the
classification table in Section~2. The distance export principle
$\delta \leq \Phi(\Gamma)$ from Definition~2.3 is concretely realized
through the AQFT framework, and the exported metric $D(\mathcal{B}_{\rm
AQFT})$ from Theorem~2.5 coincides with the HD correlation space.

\subsubsection{The AQFT Base Structure}\label{subsec:base-structure-aqft}

\begin{definition}[AQFT-based base structure for PTA]\label{def:base-aqft}
The gravitational-wave discrimination base structure
$\mathcal{B}_{\rm AQFT} = (\mathbf{C}_{\rm AQFT}, \mathcal{F}_{\rm AQFT},
\mathcal{O}_{\rm AQFT})$ is defined by:
\begin{enumerate}
\item \textbf{Source category:}
$\mathbf{C}_{\rm AQFT} = \mathbf{Loc}_M^{\mathrm{ext}}$ is the extended
category of causal regions in the PTA spacetime manifold $M$, with
morphisms given by isometric embeddings. Objects include submanifolds of
all codimensions (worldlines, world-sheets, and spacetime regions).

\item \textbf{Structure functor:}
$\mathcal{F}_{\rm AQFT} = F: \mathbf{Loc}_M^{\mathrm{ext}} \to
\mathbf{C^*Alg}$ is the extended AQFT functor (Brunetti, Fredenhagen \&
Verch~\citeyear{Brunetti:2003}) assigning to each region $O \subseteq M$
the $C^*$-algebra $\mathcal{A}(O)$ of observables localized in $O$,
satisfying Einstein causality: if $O_1, O_2$ are spacelike separated,
$[\mathcal{A}(O_1), \mathcal{A}(O_2)] = 0$.

\item \textbf{Observation functor:} For a PTA with $N$ pulsars, let
$w_i: \mathbb{R} \to M$ be the worldline of pulsar $i$. The observation
functor maps detector pairs to timing-residual correlations:
\begin{equation}
\mathcal{O}_{\rm AQFT}(w_i, w_j) = \langle R_i R_j \rangle \in \mathbb{R}
\end{equation}
where $R_i(t) \in F(w_i \cap I_t^+)$ is the timing residual at pulsar
$i$ and $I_t^+$ is the future lightcone from the Earth--pulsar baseline
crossing event. For $N$ pulsars,
$\mathcal{O}_{\rm AQFT}$ maps $\binom{N}{2}$ detector pairs to
$\binom{N}{2}$ real numbers (the correlation vector).
\end{enumerate}
\end{definition}

\begin{figure}[t]
\centering
\includegraphics[width=0.80\textwidth]{../D1-grav-waves/papers/2026-GW-BacQFT/../../Fig1-Hellings-Downs-curve.png}
\caption{The Hellings--Downs correlation function $\zeta(\theta)$ as the
observation distance $\delta$ in $\mathcal{B}_{\rm AQFT}$.}
\label{fig:hd-curve}
\end{figure}

\textbf{Pre-distance and observation distance:} The source pre-distance
on $\mathbf{C}_{\rm AQFT}$ is the total variation distance between
algebraic states on the local algebra of the PTA detector network:
\begin{equation}
\Gamma((O_1, O_2)) = d_{\mathrm{TV}}(\omega|_{\mathcal{A}(O_1)},
\omega|_{\mathcal{A}(O_2)}) = \sup_{A \in \mathcal{A}(O_1),
\,0 \leq A \leq 1} |\omega_1(A) - \omega_2(A)|
\end{equation}
where $\omega_1, \omega_2$ are algebraic states (vacuum states
corresponding to different GW source classes). The observation distance
on $\mathcal{O}_{\rm AQFT}(\mathbf{C}_{\rm AQFT}) \subseteq \mathbb{R}^P$
($P = \binom{N}{2}$) is the Hellings--Downs function:
\begin{equation}
\delta((w_i, w_j), (w_k, w_l)) = \zeta(\theta_{ij})
\end{equation}
where $\theta_{ij}$ is the angular separation of pulsars $i$ and $j$.

The natural transformation $\eta: \delta \Rightarrow \Phi(\Gamma)$
expresses the distance export principle (Definition~2.3) for this
example. Its construction follows from the functoriality of the AQFT
assignment:

\begin{proposition}[HD correlation as natural transformation]\label{prop:hd-nat}
Let $\mathbf{Emb}(M)$ be the category of spacetime embeddings (objects:
$(M, g)$, morphisms: isometric embeddings $\phi: M \to M'$). The
assignment $M \mapsto \zeta_M$, where $\zeta_M$ is the correlation
function on detector worldlines in spacetime $M$, defines a functor
$\mathsf{Z}: \mathbf{Emb}(M) \to \mathbf{Set}$. For any isometric
embedding $\phi: M \to M'$, the naturality square
\begin{equation}
\begin{tikzcd}
\zeta_M(w_i, w_j) \arrow[r, "\eta_\phi"] \arrow[d] & \zeta_{M'}(\phi(w_i), \phi(w_j)) \arrow[d] \\
\text{state}(\omega_M) \arrow[r, "\phi^*"] & \text{state}(\omega_{M'})
\end{tikzcd}
\end{equation}
commutes, making $\eta_\phi$ a natural transformation. In Minkowski
spacetime with plane-wave GWs, $\zeta_M$ reduces to the standard
Hellings--Downs function $\zeta(\theta_{12})$.

\textbf{Proof:} The proof is the same as Proposition 2 of the
grav-waves paper~\cite{D1-grav}. The GW metric perturbation in
transverse-traceless gauge is
\begin{equation}
h_{ij}(t, \vec{x}) = \int df \int d\Omega \sum_{A \in \{+, \times\}}
\tilde{h}_A(f, \hat{n})\, e^{2\pi i f(t - \hat{n}\cdot\vec{x}/c)}
\,e_{ij}^A(\hat{n})
\end{equation}
where $\langle \tilde{h}_A(f, \hat{n}) \tilde{h}_{A'}^*(f', \hat{n}')
\rangle = \tfrac{1}{2} \delta(f-f') \delta^{(2)}(\hat{n}, \hat{n}')
\delta_{AA'} S_h(f)$. The timing residual at pulsar $i$ is
\begin{equation}
R_i(t) = \int dt'\, h_{ij}(t', \vec{x}_i)\, \hat{n}_i^i \hat{n}_i^j
\,g(t')
\end{equation}
with response function $g(t')$. The cross-correlation for two pulsars
with angular separation $\theta_{12}$ is
\begin{equation}
\langle R_1 R_2 \rangle = \int df\, S_h(f)\, \zeta(\theta_{12}, f)
\end{equation}
In the long-wavelength limit (PTA regime, $f \to 0$), this reduces to
\begin{equation}
\zeta(\theta_{12}) = \frac{3}{2} x \ln x - \frac{x}{4} + \frac{1}{2} +
\frac{1}{2}\delta_{12}
\end{equation}
where $x = (1 - \cos\theta_{12})/2$. The $SO(3)$ rotational symmetry of
$\zeta(\theta)$ reflects the $SO(3)$ rotation group structure of the
AQFT assignment on the sphere at infinity; the logarithmic term arises
from the non-trivial cobordism class connecting the two worldlines in
$\mathbf{Loc}_M^{\mathrm{ext}}$. Functoriality of $F$ guarantees that
$\zeta$ transforms naturally under isometric embeddings, establishing
the naturality of $\eta$. $\square$
\end{proposition}

\subsubsection{Source Discrimination via Total Variation Distance}
\label{subsec:source-discrimination}

\begin{theorem}[GW source discrimination]\label{thm:source-discrim}
Let $\omega_\alpha$ be the algebraic state on the PTA's local algebra
associated with GW source class $\alpha \in \{\mathrm{SMBHB},
\mathrm{cosmo}\}$. The total variation distance between observation
distributions satisfies:
\begin{equation}
d_{\mathrm{TV}}(\omega_{\mathrm{SMBHB}}, \omega_{\mathrm{cosmo}})
\gtrsim c \cdot P^{-1/2} \cdot \bigl| \sigma_N^2(\mathrm{SMBHB}) -
\sigma_N^2(\mathrm{cosmo}) \bigr|
\end{equation}
where $P = \binom{N}{2}$ is the number of pulsar pairs, $\sigma_N^2$ is
the fractional variance of the source count in the PTA sky coverage, and
$c$ is a constant depending on the spectral-index difference
$\Delta\gamma = |\gamma_{\mathrm{SMBHB}} - \gamma_{\mathrm{cosmo}}|$.

The discrimination threshold follows from the inequality
\begin{equation}
P^{-1/2} \lesssim \bigl| \sigma_N^2(\mathrm{SMBHB}) -
\sigma_N^2(\mathrm{cosmo}) \bigr|
\end{equation}
which is a direct consequence of the distance export principle: the
observable distance $\delta = \zeta(\theta_{12})$ can distinguish two
source classes only when their source pre-distances $\Gamma$ produce a
difference exceeding the statistical noise floor $P^{-1/2}$.
\end{theorem}

\begin{proof}
By Le Cam's lemma, $d_{\mathrm{TV}}$ characterizes the minimum error
probability in binary hypothesis testing. For the HD correlation with
$P$ pulsar pairs, the statistical uncertainty is $\delta\zeta \sim
P^{-1/2}$. Distinguishing two source classes requires their HD
correlation functions to differ by more than this noise floor:
\begin{equation}
\delta\zeta \lesssim \bigl| \zeta_{\mathrm{SMBHB}}(\theta) -
\zeta_{\mathrm{cosmo}}(\theta) \bigr| \sim
\bigl| \sigma_N^2(\mathrm{SMBHB}) - \sigma_N^2(\mathrm{cosmo}) \bigr|
\end{equation}

\textbf{SMBHBs.} Sources are discrete ($N_{\mathrm{src}} \sim
10^4\text{--}10^6$ at nanohertz frequencies) located along the galaxy
distribution. The angular distribution traces large-scale structure,
producing Poisson shot noise with fractional variance
$\sigma_N^2(\mathrm{SMBHB}) \sim \mathcal{O}(N_{\mathrm{src}}^{-1})$.

\textbf{Cosmological backgrounds.} The anisotropy is dominated by the
projection of the primordial density field onto the sky. For
inflationary GWBs, the angular power spectrum $C_\ell \sim
\mathcal{O}(10^{-10})$ for $\ell \lesssim 100$ (PTA-relevant
multipoles), set by the Sachs--Wolfe effect. The absolute anisotropy
amplitude is determined by primordial density fluctuations
$\delta\rho/\rho \sim \mathcal{O}(10^{-5})$, independent of source
count.

By the distance export principle (Definition~2.3), the natural
transformation $\eta$ relates the observable HD correlation $\delta$ to
the source distance $\Gamma = d_{\mathrm{TV}}$. When
$d_{\mathrm{TV}}(\omega_{\mathrm{SMBHB}}, \omega_{\mathrm{cosmo}}) >
P^{-1/2}$, the exported metric $D(\mathcal{B}_{\mathrm{AQFT}})$
(Theorem~2.5) yields distinguishable observations. When the inequality
reverses, Theorem~2.6 (observation infeasibility) applies: distinct base
structures with different source classes yield isometric observation
spaces, and observation alone cannot discriminate them. $\square$
\end{proof}

\begin{remark}[Relation to DistCat observation infeasibility]
Theorem~2.6 states that there exist $\mathcal{B}_1 \neq \mathcal{B}_2$
with $D(\mathcal{B}_1) \cong D(\mathcal{B}_2)$. In this example, when
$P^{-1/2} > |\sigma_N^2(\mathrm{SMBHB}) - \sigma_N^2(\mathrm{cosmo})|$,
the HD correlation vectors for the two source classes are observationally
indistinguishable despite their different algebraic origins. This is the
physical realization of observation infeasibility.
\end{remark}

\subsubsection{Isotropy Scaling Theorem}\label{subsec:isotropy-scaling}

\begin{theorem}[Isotropy scaling for cosmological GWB]\label{thm:isotropy-scaling}
A cosmological GWB (from inflation, phase transitions, or other
early-universe processes) produces an HD correlation with anisotropic
corrections determined by the large-scale structure projection:
\begin{equation}
\zeta_{\rm cosmo}(\theta_{12}) = \zeta_0(\theta_{12}) +
\mathcal{O}\!\bigl(C_\ell^{\mathrm{cosmo}}\bigr) \cdot
\mathcal{Y}_{\ell m}^{\mathrm{cosmo}}(\hat{n}_{12})
\end{equation}
where $\zeta_0$ is the standard isotropic HD function,
$C_\ell^{\mathrm{cosmo}}$ is the angular power spectrum of the
cosmological anisotropy, and $\mathcal{Y}_{\ell m}^{\mathrm{cosmo}}$
are the corresponding spherical-harmonic coefficients.

For an SMBHB GWB:
\begin{equation}
\zeta_{\mathrm{SMBHB}}(\theta_{12}) = \zeta_0(\theta_{12}) +
\mathcal{O}(N_{\mathrm{src}}^{-1/2}) \cdot
\mathcal{Y}_{\ell m}^{\mathrm{LSS}}(\hat{n}_{12})
\end{equation}
where $\mathcal{Y}_{\ell m}^{\mathrm{LSS}}$ traces the large-scale
structure of the local galaxy distribution.

The anisotropy amplitudes differ by four orders of magnitude:
\begin{equation}
\frac{\text{anisotropy}_{\mathrm{SMBHB}}}{\text{anisotropy}_{\mathrm{cosmo}}}
\sim \frac{\mathcal{O}(1\text{--}10\%)}{\mathcal{O}(10^{-5})} \sim
\mathcal{O}(10^3\text{--}10^4)
\end{equation}

\textbf{DistCat interpretation:} The anisotropy amplitude ratio is a
2-morphism invariant under the distance export functor. Since
$\delta_{\mathrm{SMBHB}} \neq \delta_{\mathrm{cosmo}}$ at the level of
observable distances (Theorem~\ref{thm:source-discrim}), the two source
classes belong to \textbf{distinct distance-export equivalence classes}
in $\mathbf{DistCat}$, even though they share the same base category
$\mathbf{C}_{\mathrm{AQFT}}$ and structure functor $\mathcal{F}_{\mathrm{AQFT}}$.
\end{theorem}

\begin{proof}
\textbf{Cosmological GWB.} Sources are distributed at high redshift
($z \gtrsim 10$ for inflation, $z \sim 1\text{--}10$ for phase
transitions). By the cosmological principle, the source distribution is
homogeneous and isotropic to $\mathcal{O}(10^{-5})$. The anisotropy in
$\zeta$ arises from (1) primordial density fluctuations that generated
the GW, and (2) gravitational redshift and Shapiro delay along the line
of sight. For inflationary GWBs, $C_\ell \sim \mathcal{O}(10^{-10})$
dominated by the Sachs--Wolfe effect at $\ell \lesssim 100$, yielding
fractional anisotropy $\mathcal{O}(10^{-5})$.

\textbf{SMBHBs.} Sources are galaxies in the local universe ($z \lesssim
1$), with $N_{\mathrm{src}} \sim 10^4\text{--}10^6$ in the PTA volume.
The galaxy distribution traces the cosmic web with significant
anisotropy: local supercluster structure produces
$\mathcal{O}(1\text{--}10\%)$ fluctuations on PTA-relevant angular
scales, scaling as $N_{\mathrm{src}}^{-1/2}$ (Poisson shot noise).

The key discriminant is both the amplitude and the \textbf{structure} of
the anisotropy: SMBHB anisotropy directly traces the galaxy LSS with
$\mathcal{O}(1\text{--}10\%)$ amplitude and correlates with galaxy
surveys, while cosmological anisotropy is set by the primordial density
field at $\mathcal{O}(10^{-5})$.

\textbf{DistCat proof:} In the DistCat framework, the anisotropy is a
property of the observation distance $\delta = \zeta(\theta_{12})$. The
2-functor $D: \mathbf{DistCat} \to \mathbf{Met}$ (Theorem~2.5) maps
$\mathcal{B}_{\mathrm{SMBHB}}$ and $\mathcal{B}_{\mathrm{cosmo}}$ to
metric spaces whose observation distances differ by the anisotropy
amplitude. Since $D(\mathcal{B}_{\mathrm{SMBHB}}) \not\cong
D(\mathcal{B}_{\mathrm{cosmo}})$ as metric spaces (the distance
functions are distinct functions on the observation space), the two
base structures are non-isomorphic in the metric-exported category. The
anisotropy amplitude ratio is thus a DistCat invariant. $\square$
\end{proof}

\begin{corollary}[Observability of SMBHB anisotropy]\label{cor:skao-obs}
For NANOGrav 15-year data ($N = 68$, $P = 2278$), the discrimination
SNR for SMBHB vs cosmological anisotropy is
\begin{equation}
\mathrm{SNR}_{\mathrm{dis}} \sim P^{1/2} \cdot \mathcal{O}(1\text{--}10\%)
\sim 30\text{--}300
\end{equation}
For SKAO-era observations ($N \sim 200$, $P \sim 20000$):
\begin{equation}
\mathrm{SNR}_{\mathrm{dis}} \sim \begin{cases}
\sim 140 & (N_{\mathrm{src}} \sim 10^4)\\
\sim 45 & (N_{\mathrm{src}} \sim 10^5)\\
\sim 14 & (N_{\mathrm{src}} \sim 10^6)
\end{cases}
\end{equation}
Current data sensitivity to HD anisotropy is $\sim 0.1$
\cite{Isi:2024}, so SMBHB anisotropy may be detectable with next-
generation data if systematics are controlled.
\end{corollary}

\subsubsection{Summary and DistCat Properties}\label{subsec:summary}

Table~\ref{tab:grav-distcat} summarizes the DistCat instantiation for
gravitational-wave source discrimination.

{\def\LTcaptype{none}
\begin{longtable}{@{}
  >{\raggedright\arraybackslash}p{0.28\linewidth}
  >{\raggedright\arraybackslash}p{0.70\linewidth}@{}}
\toprule
\textbf{DistCat component} & \textbf{Grav-waves instantiation} \\
\midrule
\endhead
\bottomrule
\endlastfoot
Base structure $\mathcal{B}_{\rm AQFT}$ &
$(\mathbf{Loc}_M^{\mathrm{ext}},\; F,\; \langle R_i R_j \rangle)$ \\
Source category $\mathbf{C}$ &
Extended causal regions $\mathbf{Loc}_M^{\mathrm{ext}}$ \\
Structure functor $\mathcal{F}$ &
Extended AQFT: $O \mapsto \mathcal{A}(O) \in \mathbf{C^*Alg}$ \\
Observation functor $\mathcal{O}$ &
$\mathcal{O}(w_i, w_j) = \langle R_i R_j \rangle = \zeta(\theta_{ij})$ \\
Source pre-distance $\Gamma$ &
$d_{\mathrm{TV}}(\omega_1, \omega_2)$ on $\mathcal{A}(\text{PTA})$ \\
Observation distance $\delta$ &
$\zeta(\theta_{12})$ (Hellings--Downs correlation) \\
Natural transformation $\eta$ &
$\eta: \zeta \Rightarrow \Phi(d_{\mathrm{TV}})$ (local covariance) \\
Exported metric $D(\mathcal{B})$ &
$(\{\zeta(\theta_{ij})\}, \delta_{\mathrm{HD}})$ \\
Distance equivalence class &
Distinguished by anisotropy amplitude (Thm~\ref{thm:isotropy-scaling}) \\
Observation infeasibility &
$D(\mathcal{B}_{\mathrm{SMBHB}}) \cong D(\mathcal{B}_{\mathrm{cosmo}})$
when $d_{\mathrm{TV}} < P^{-1/2}$ (Thm~2.6) \\
\end{longtable}
}

The five examples in Section~3 establish $\delta \sim \Gamma^\alpha$ with
critical exponents $\alpha \in \{1.84\}$ for this example (log-corrected
quadratic). The universal pattern $\delta \sim \Gamma^\alpha$ across
all examples (Section~8) suggests that the distance export principle is
a categorical invariant.

\begin{thebibliography}{99}

\bibitem{D1-grav}
C.~Lu. Extended functorial QFT framework for GW source discrimination in
pulsar timing arrays. \textit{in preparation} (D1-grav-waves).

\end{thebibliography}

